\section{Results}

This section reports the correctness, transition, topology, and token-cost results for the four Research Questions.
Each subsection begins with the aggregate pattern and then examines variation or secondary evidence needed to interpret that pattern.

\subsection{Stage Composition and Regression Control}

Table~\ref{tab:rq1-stage-results} summarizes two complementary views of the Separate-Call protocols.
Panel A compares each outcome with its Shared Initial Candidate, whereas Panel B reports paired differences between adjacent Stage Compositions.
The pooled rows provide overall magnitudes; the model-benchmark comparisons remain the primary evidence of heterogeneity.

\begin{table}[H]
\caption{Pooled correctness and transition results for Separate-Call Stage Composition.}
\label{tab:rq1-stage-results}
\centering
\begingroup
\scriptsize
\renewcommand{\arraystretch}{1.12}
\begin{threeparttable}
\textit{Panel A: Outcomes relative to the Shared Initial Candidate}\\[2pt]
\begin{tabularx}{0.94\columnwidth}{
    >{\raggedright\arraybackslash}p{0.14\columnwidth}
    >{\centering\arraybackslash}p{0.15\columnwidth}
    >{\centering\arraybackslash}p{0.13\columnwidth}
    >{\centering\arraybackslash}X
    >{\centering\arraybackslash}X
    >{\centering\arraybackslash}X
}
\toprule
Outcome & Determinate & Pass rate & Repairs & Regressions & Repair $-$ regression \\
\midrule
Direct & 10,033 & 56.44\% & 0 & 0 & 0 \\
\texttt{R} & 10,032 & 56.79\% & 56 & 22 & +34 \\
\texttt{CR} & 10,029 & 52.54\% & 222 & 610 & $-$388 \\
\texttt{CPR} & 10,027 & 51.31\% & 209 & 725 & $-$516 \\
\bottomrule
\end{tabularx}

\vspace{5pt}
\textit{Panel B: Paired differences between adjacent compositions}\\[2pt]
\begin{tabularx}{0.94\columnwidth}{
    >{\raggedright\arraybackslash}p{0.20\columnwidth}
    >{\centering\arraybackslash}p{0.14\columnwidth}
    >{\centering\arraybackslash}X
    >{\centering\arraybackslash}p{0.14\columnwidth}
    >{\centering\arraybackslash}p{0.14\columnwidth}
}
\toprule
Comparison & Paired $N$ & $\Delta$ correct, pp [95\% CI] & Improved & Worsened \\
\midrule
Direct to \texttt{R} & 10,032 & +0.34 [0.16, 0.52] & 56 & 22 \\
\texttt{R} to \texttt{CR} & 10,021 & $-$4.26 [$-$4.81, $-$3.71] & 188 & 615 \\
\texttt{CR} to \texttt{CPR} & 10,023 & $-$1.24 [$-$1.61, $-$0.87] & 112 & 236 \\
\bottomrule
\end{tabularx}
\begin{tablenotes}[flushleft]
\footnotesize
\item Pass rates exclude indeterminate functional outcomes.
Repairs and Regressions in Panel A are transitions from the Initial Candidate and use complete-transition denominators.
Improved and Worsened in Panel B instead compare the two named outcomes on their paired-complete denominator.
Confidence intervals are paired bootstrap intervals.
\end{tablenotes}
\end{threeparttable}
\endgroup
\end{table}

Revision alone produced the highest pooled determinate pass rate among the Separate-Call refinement protocols.
Relative to Direct Generation, \texttt{R} increased correctness by 0.34 percentage points, with 56 improved and 22 worsened pairs; the exact McNemar test also rejected equality ($p=0.00015$).
The magnitude was nevertheless small and heterogeneous: the difference was positive in 11 of the 18 model-benchmark combinations, negative in three, and zero in four, and only one individual combination had $p<0.05$.

Adding Critique before Revision reversed the net result.
The \texttt{CR} protocol was 4.26 percentage points below \texttt{R} and was negative in 17 of 18 model-benchmark combinations.
Adding Planning reduced correctness by a further 1.24 percentage points, with \texttt{CPR} below \texttt{CR} in 14 combinations.
These declines did not reflect an inability to repair code: \texttt{CR} and \texttt{CPR} repaired 222 and 209 initially incorrect candidates, compared with 56 for \texttt{R}.
Instead, they regressed 610 and 725 initially correct candidates, compared with 22 for \texttt{R}.
The central RQ1 result is therefore that regression control, rather than the number of repairs alone, determined whether added stages improved final correctness.

This finding aligns with the transition-based perspective of CoCoS, which evaluates correction together with damage to correct outputs \cite{cho-etal-2025-self}, and with CodeChat-Eval's observation that repeated functionality-preserving edits can reduce Functional Correctness \cite{guo2026codechateval}.
Our result establishes the same preservation concern in a different setting: no correction behavior was trained, the Critique and Plan were self-generated without Execution Feedback, and every protocol revised the same Initial Candidate.
The comparisons are complete protocol effects and do not establish whether a harmful outcome originated in an upstream artifact or in Revision's use of that artifact.

\subsection{Decision Benefits Regression-Prone Refinement}

Table~\ref{tab:rq2-decision-results} compares each Always-Refine protocol with the outcome obtained by applying the shared Decision.
Prevented Regression counts an initially correct candidate that Always-Refine damaged but Decision preserved, whereas Missed Repair counts an initially incorrect candidate that Always-Refine repaired but Decision preserved.
Token savings represent the downstream calls omitted after \texttt{PRESERVE} under protocol-implied execution.

\begin{table}[H]
\caption{Correctness and cost effects of Decision conditioning.}
\label{tab:rq2-decision-results}
\centering
\begingroup
\scriptsize
\renewcommand{\arraystretch}{1.12}
\begin{threeparttable}
\begin{tabularx}{0.98\columnwidth}{
    >{\raggedright\arraybackslash}p{0.16\columnwidth}
    >{\centering\arraybackslash}X
    >{\centering\arraybackslash}p{0.13\columnwidth}
    >{\centering\arraybackslash}p{0.12\columnwidth}
    >{\centering\arraybackslash}p{0.13\columnwidth}
    >{\centering\arraybackslash}p{0.13\columnwidth}
}
\toprule
Comparison & $\Delta$ correct, pp [95\% CI] & Prevented Regression & Missed Repair & Tokens saved, million & Mean saved \\
\midrule
\texttt{R} to \texttt{DR} & $-$0.13 [$-$0.27, 0.00] & 18 & 31 & 1.93 & 192 \\
\texttt{CR} to \texttt{DCR} & +3.95 [3.45, 4.45] & 540 & 149 & 11.59 & 1,153 \\
\texttt{CPR} to \texttt{DCPR} & +5.10 [4.55, 5.65] & 646 & 136 & 18.84 & 1,874 \\
\bottomrule
\end{tabularx}
\begin{tablenotes}[flushleft]
\footnotesize
\item Correctness differences use paired-complete outcomes.
Transition counts use the corresponding complete Decision and Initial-to-final outcomes and can differ slightly from the paired correctness counts.
Mean savings are protocol-implied tokens per model-task unit.
\end{tablenotes}
\end{threeparttable}
\endgroup
\end{table}

Decision was strongly conservative: 9,244 of 10,054 resolved labels were \texttt{PRESERVE} (91.94\%), and 810 were \texttt{REFINE}.
For the relatively stable \texttt{R} path, this policy prevented 18 Regressions but missed 31 Repairs, resulting in a small difference whose confidence interval included zero.
It nevertheless saved 1.93 million protocol-implied tokens.

The same Decision was beneficial when applied to the more regression-prone paths.
\texttt{DCR} recovered 3.95 percentage points over \texttt{CR} while saving 11.59 million tokens, and \texttt{DCPR} recovered 5.10 points over \texttt{CPR} while saving 18.84 million.
The \texttt{DCR} difference was positive in 17 of 18 model-benchmark combinations, and the \texttt{DCPR} difference was positive in all 18.
Restricting the analysis to Decisions accepted by the exact parser produced nearly identical differences of $-$0.13, +4.01, and +5.15 percentage points for \texttt{DR}, \texttt{DCR}, and \texttt{DCPR}, respectively.

Initial Pass Rate alone did not explain this pattern.
Across the 18 combinations, its descriptive Pearson correlations with the \texttt{DCR} and \texttt{DCPR} correctness gains were $-$0.13 and $-$0.14.
Its correlations with mean token savings were moderately negative for both paths ($r=-0.53$), contrary to a simple account in which more initially correct candidates necessarily imply larger savings.
These associations are descriptive and do not separate model, benchmark, tokenization, and Decision behavior.
The stronger empirical predictor in the protocol comparisons was the downstream path's observed Regression risk.

The coexistence of prevented Regressions and missed Repairs aligns with prior evidence that execution-free LLM judgments can reject correct implementations for unsupported reasons \cite{jin2026reliable}.
That study evaluates review verdicts, whereas our Shared-Initial Design follows an independent binary Decision through counterfactual candidate selection.
The present result therefore extends the preservation concern to final code outcomes while also showing that an imperfect Decision can be useful when the alternative refinement path is sufficiently destructive.

\subsection{Single-Call and Separate-Call Topology}

Table~\ref{tab:rq3-topology-results} reports matched topology comparisons and the consistency of exact Decision labels with emitted code behavior.
Panel A compares final correctness under the Separate-Call and Single-Call realizations, while Panel B records whether the code changed in the direction implied by each Single-Call Decision label.

\begin{table}[H]
\caption{Correctness effects of call topology and consistency of Single-Call Decision labels.}
\label{tab:rq3-topology-results}
\centering
\begingroup
\scriptsize
\renewcommand{\arraystretch}{1.12}
\begin{threeparttable}
\textit{Panel A: Single-Call minus matched Separate-Call correctness}\\[2pt]
\begin{tabularx}{0.94\columnwidth}{
    >{\raggedright\arraybackslash}p{0.23\columnwidth}
    >{\centering\arraybackslash}p{0.15\columnwidth}
    >{\centering\arraybackslash}X
    >{\centering\arraybackslash}p{0.14\columnwidth}
    >{\centering\arraybackslash}p{0.14\columnwidth}
}
\toprule
Comparison & Paired $N$ & $\Delta$ correct, pp [95\% CI] & Improved & Worsened \\
\midrule
\texttt{CR} to \texttt{SC-CR} & 10,023 & +3.95 [3.41, 4.49] & 607 & 211 \\
\texttt{CPR} to \texttt{SC-CPR} & 10,020 & +3.21 [2.58, 3.83] & 667 & 345 \\
\texttt{DR} to \texttt{SC-DR} & 10,032 & +0.07 [$-$0.11, 0.25] & 43 & 36 \\
\texttt{DCR} to \texttt{SC-DCR} & 10,029 & +0.05 [$-$0.22, 0.33] & 100 & 95 \\
\texttt{DCPR} to \texttt{SC-DCPR} & 10,027 & $-$0.09 [$-$0.40, 0.22] & 114 & 123 \\
\bottomrule
\end{tabularx}

\vspace{5pt}
\textit{Panel B: Exact Decision-label consistency with emitted code}\\[2pt]
\begin{tabularx}{0.94\columnwidth}{
    >{\raggedright\arraybackslash}p{0.18\columnwidth}
    >{\centering\arraybackslash}p{0.16\columnwidth}
    >{\centering\arraybackslash}X
    >{\centering\arraybackslash}X
    >{\centering\arraybackslash}p{0.17\columnwidth}
}
\toprule
Protocol & Exact labels & \texttt{PRESERVE} with changed code & \texttt{REFINE} with unchanged code & Inconsistency rate \\
\midrule
\texttt{SC-DR} & 10,054 & 180 & 674 & 8.49\% \\
\texttt{SC-DCR} & 10,054 & 123 & 763 & 8.81\% \\
\texttt{SC-DCPR} & 10,054 & 114 & 297 & 4.09\% \\
\bottomrule
\end{tabularx}
\begin{tablenotes}[flushleft]
\footnotesize
\item Positive differences in Panel A favor Single-Call.
Label inconsistency is behavioral: it compares the exact label with source-code equality to the Initial Candidate and does not evaluate whether changing the code was semantically justified.
\end{tablenotes}
\end{threeparttable}
\endgroup
\end{table}

Combining Critique and Revision in one call recovered 3.95 percentage points relative to \texttt{CR}, and combining Critique, Planning, and Revision recovered 3.21 points relative to \texttt{CPR}.
The first difference was positive in 17 of 18 model-benchmark combinations, and the second was positive in 16 and negative in two.
This pattern is consistent with avoiding some failure introduced by externalized artifact chains, but the experiment does not isolate artifact quality from prompt context, call boundaries, or Revision behavior.
It also does not establish that the model internally performed Critique or Planning in the requested order.

No aggregate correctness advantage appeared once Decision was part of the matched protocols.
The pooled differences for \texttt{SC-DR}, \texttt{SC-DCR}, and \texttt{SC-DCPR} were all within 0.1 percentage points of their Separate-Call counterparts, and every confidence interval included zero.
Single-Call labels were not fully aligned with emitted source changes: behavior contradicted the exact label in 8.49\% of \texttt{SC-DR}, 8.81\% of \texttt{SC-DCR}, and 4.09\% of \texttt{SC-DCPR} responses.
Applying the labels after generation changed the number of correct outcomes by only +2, +5, and $-$2 on the corresponding complete sensitivity subsets.
Thus, the separate Decision call did not improve pooled correctness but retained an auditable exact-selection boundary that the combined response did not guarantee behaviorally.

The absence of a universal topology winner is consistent with Feedback Over Form, which found no clear advantage for more elaborate loop topology in its local-model experiments \cite{mcandrews2026feedback}.
Its experiments used models with 1B to 3B parameters and Execution Feedback, whereas ours used larger open-weight models and no behavioral feedback.
Within our setting, call topology mattered specifically for the Critique and Planning artifact chains, but not for Decision-bearing protocols in aggregate.

\subsection{Correctness Relative to Token Cost}

The final physical experiment consumed 91,816,683 tokens, comprising 65,295,167 input and 26,521,516 output tokens.
Table~\ref{tab:rq4-cost-results} compares protocol-implied correctness and token use on the 10,006 model-task units with determinate outcomes for every protocol.
The Pareto column marks protocols for which no other protocol was both at least as correct and no more expensive, with one strict advantage.

\begin{table}[H]
\caption{Pooled correctness-cost comparison on the common-complete subset.}
\label{tab:rq4-cost-results}
\centering
\begingroup
\scriptsize
\renewcommand{\arraystretch}{1.08}
\begin{threeparttable}
\begin{tabularx}{0.76\columnwidth}{
    >{\raggedright\arraybackslash}X
    >{\centering\arraybackslash}p{0.18\columnwidth}
    >{\centering\arraybackslash}p{0.23\columnwidth}
    >{\centering\arraybackslash}p{0.14\columnwidth}
}
\toprule
Protocol & Correct & Tokens per model-task & Pareto \\
\midrule
Direct & 56.58\% & 460 & Yes \\
\texttt{DR} & 56.79\% & 974 & Yes \\
\texttt{DCR} & 56.61\% & 1,098 & No \\
\texttt{R} & 56.92\% & 1,166 & Yes \\
\texttt{DCPR} & 56.52\% & 1,186 & No \\
\texttt{SC-CR} & 56.61\% & 1,200 & No \\
\texttt{SC-DR} & 56.86\% & 1,217 & No \\
\texttt{SC-DCR} & 56.66\% & 1,234 & No \\
\texttt{SC-CPR} & 54.63\% & 1,254 & No \\
\texttt{SC-DCPR} & 56.43\% & 1,255 & No \\
\texttt{CR} & 52.65\% & 2,250 & No \\
\texttt{CPR} & 51.41\% & 3,059 & No \\
\bottomrule
\end{tabularx}
\begin{tablenotes}[flushleft]
\footnotesize
\item Direct includes Initial Candidate generation.
Every other row includes Direct plus the calls implied by that protocol.
Tokens are model-tokenizer counts and are not hardware-normalized compute measures.
\end{tablenotes}
\end{threeparttable}
\endgroup
\end{table}

Only Direct, \texttt{DR}, and \texttt{R} lay on the pooled correctness-cost frontier.
Direct was the low-cost anchor, \texttt{R} had the highest correctness, and \texttt{DR} provided an intermediate cost and correctness point.
Every protocol containing an external Critique or Plan was dominated in the pooled common-complete comparison.
The Single-Call protocols were much cheaper than their corresponding multi-call \texttt{CR} and \texttt{CPR} paths, but none exceeded the frontier established by Direct, \texttt{DR}, and \texttt{R}.

Relative to Direct, \texttt{R} added 34 correct outcomes at an additional 7.06 million tokens, or about 208 thousand tokens per additional correct outcome.
The corresponding \texttt{DR} increment was 21 correct outcomes and 5.14 million tokens, or about 245 thousand tokens per additional correct outcome.
These ratios describe pooled token accounting rather than monetary or hardware-normalized efficiency, and a deployment choice can differ when the cost of an incorrect solution is not linear in token count.

The limited return from additional refinement computation aligns with prior evidence that repeated self-repair is not uniformly more effective than alternative use of a generation budget \cite{olausson2024selfrepair}.
That work compares repair with independent resampling and often uses execution-derived feedback, whereas our comparison holds one Initial Candidate fixed, supplies no Execution Feedback, and accounts for the calls in each Stage Composition.
Within this narrower setting, additional Critique and Planning calls did not justify their token cost in the pooled panel, even though Single-Call composition improved correctness relative to its matched multi-call path.
